% %--------------------------
% %	CONFIGURATION
% %--------------------------
\documentclass[11pt,a4paper]{moderncv}
\moderncvstyle{classic}
\moderncvcolor{black}
\definecolor{firstnamecolor}{rgb}{0,0,0}

% Basic packages
\usepackage[utf8]{inputenc}
\usepackage[T1]{fontenc}
\usepackage[scale=0.85]{geometry}

\usepackage{microtype}

% Configure font settings
\renewcommand{\rmdefault}{ppl}
\renewcommand{\sfdefault}{phv}
\renewcommand*{\namefont}{\fontsize{26}{18}\mdseries}

% Configure microtype
\microtypesetup{expansion=false}

% Load country-specific contact information after copying into outputs/<listing_slug>/.
% Use personal_info_dk.tex for Denmark-facing applications and personal_info_se.tex for Sweden-facing applications.
\input{../../secrets/personal_info_dk.tex}

% Title document
\title{Curriculum Vitae}

% %--------------------------

\begin{document}
\makecvtitle

\section*{Introduction}
Quantitative researcher with a background in Engineering Physics and Statistical Learning and Data Analytics from KTH, plus an exchange in Applied Mathematics at ETH. Since 2019, I have worked on machine learning in medical imaging, pharmaceutical analytics, and production systems. I am especially interested in mathematically grounded image analysis, quantitative characterisation, and benchmarking in scientific settings.

\section{Education}

\cventry{2015--2020}{Royal Institute of Technology, KTH -- Engineering Physics}{}{GPA: 5.0 (out of 5.0)}{}{Master of Science: Statistical Learning and Data Analytics; Bachelor of Science: Engineering Physics. \href{https://kth.diva-portal.org/smash/get/diva2:1431327/FULLTEXT02.pdf}{See thesis}. \newline Degree project: \emph{Image Distance Learning for Probabilistic Dose-Volume Histogram and Spatial Dose Prediction in Radiation Therapy Treatment Planning}. \newline Relevant coursework: Statistical Machine Learning; Modern Methods of Statistical Learning; Computational Statistics.}

\cventry{2019}{Swiss Federal Institute of Technology, ETH -- Applied Mathematics}{Exchange}{}{}{}

\cventry{2017--2022}{Stockholm University, SU -- Economics}{Supplementary Bachelor's Degree}{}{}{}

% ---------------------------
% Selected Work Experience
% ---------------------------
\section{Selected Work Experience}

% ---------------------------
\large{Researcher within Data Science @ RaySearch Laboratories}
\normalsize
\begin{itemize}
\item Worked on automated radiotherapy planning using deep learning, probabilistic modelling, and multi-objective optimisation.
\item Built autoencoder-based pipelines for latent anatomical features from segmented 3D CT scans and learned similarity metrics between patients.
\item Predicted dose-volume histograms and spatial dose distributions with sparse Gaussian processes.
\end{itemize}

{\footnotesize\textit{Keywords: medical imaging, autoencoders, sparse Gaussian processes, optimisation, clinical collaboration}}

\hfill{\small{\textit{\hyperref[sec:raysearch]{Read more...}}}}

% ---------------------------
\large{Full-Stack Data Scientist \& Project Manager @ Signum Life Science}
\normalsize
\begin{itemize}
\item Maintain and develop a pharmaceutical forecasting platform covering price, demand, and patient population forecasts.
\item Use LLM embeddings for downstream similarity search and matching workflows.
\item Support commercialisation and internal collaboration through technical ownership and project management.
\end{itemize}

{\footnotesize\textit{Keywords: production ML, embeddings, similarity search, forecasting, Databricks, AWS}}

\hfill{\small{\textit{\hyperref[sec:signum]{Read more...}}}}

% ---------------------------
\large{Principal Data Scientist @ Valcon}
\normalsize
\begin{itemize}
\item Led data science projects across pharma, energy, IoT, and manufacturing.
\item Helped build an end-to-end production pipeline for a real-time deep neural network.
\item Mentored junior colleagues, led explainable AI and Git training, and presented technical work to non-technical stakeholders.
\end{itemize}

{\footnotesize\textit{Keywords: production ML, explainable AI, mentoring, stakeholder communication, deep neural networks}}

\hfill{\small{\textit{\hyperref[sec:valcon]{Read more...}}}}

% ---------------------------
% Skills and Other
% ---------------------------
\section{Skills and Other}
\normalsize
\cvitem{Technology}{Python, SQL, PyTorch, TensorFlow, Scikit-learn, XGBoost, SHAP, autoencoders, CNNs, sparse Gaussian processes, Databricks, AWS, FastAPI, Docker, Qdrant}
\cvitem{Languages}{\textit{Native}: English, Swedish. \textit{Intermediate}: Danish.}
\cvitem{Other}{Valcon People's Choice Award; member of Mensa Sweden.}

\newpage

% ---------------------------
% Detailed Work Experience
% ---------------------------
\section{Detailed Work Experience}

% ---------------------------
\phantomsection\label{sec:raysearch}
\cventry{2019--2021}{Researcher within Data Science}{RaySearch Laboratories}{Stockholm}{}{
    At RaySearch Laboratories, I worked on automating radiotherapy planning using deep learning, probabilistic modelling, and multi-objective optimisation. I developed a pipeline that extracted latent anatomical features from segmented 3D CT scans using autoencoders, and used those representations to identify similar patients via a learned similarity metric. Dose-volume histograms and spatial dose distributions were then predicted using sparse Gaussian processes, providing a probabilistic framework for dose planning conditioned on patient geometry.
    \newline\hspace*{2em}
    My degree project at KTH was \emph{Image Distance Learning for Probabilistic Dose-Volume Histogram and Spatial Dose Prediction in Radiation Therapy Treatment Planning}, which aligns closely with this work. I also collaborated closely with medical physicists and clinicians to keep the models aligned with oncological standards and practical workflow requirements.
}{}

% ---------------------------
\phantomsection\label{sec:signum}
\cventry{2024--Present}{Full-Stack Data Scientist \& Project Manager}{Signum Life Science}{Copenhagen}{}{
    At Signum, I maintain and lead development on a pharmaceutical forecasting platform covering price, demand, and patient population forecasts. I work across architecture, project management, data science, ML engineering, and software engineering for the platform.
    \newline\hspace*{2em}
    In my current work, I am using LLMs to embed free text for downstream similarity search and matching workflows. I am also supporting commercialisation by helping sales representatives use the system in practice.
    \newline\hspace*{2em}
    Beyond model development, I organised an internal week-long hackathon with up to 40 participants.
}{}

% ---------------------------
\phantomsection\label{sec:valcon}
\cventry{2022--2024}{Principal Data Scientist}{Valcon}{Copenhagen}{}{
    At Valcon, I led data science projects across pharmaceuticals, energy, IoT, and manufacturing. I also mentored junior colleagues, led explainable AI and Git trainings, and regularly presented technical topics to non-technical stakeholders.
    \newline\hspace*{2em}
    One project involved training a real-time deep neural network with a custom training pipeline and modular architecture. The training data consisted of natural-language inputs and high-dimensional multi-class labels, which made for a challenging prediction problem. I was responsible for the end-to-end production pipeline.
    \newline\hspace*{2em}
    I was awarded the Valcon People's Choice Award for exemplary work, leadership, team spirit, and scientific curiosity.
}{}

\end{document}
